\begin{center}\begin{tikzpicture}[font=\small,>=stealth,line join=round]
\draw[gray!40,rounded corners=3pt] (0,0) rectangle (14.2,8.3);\node at (7.1,7.700000000000001) {例4：关于负1与3反射两次，横向总共移动8};
\begin{scope}[shift={(2.7153846153846155,3.8222222222222224)},x=0.46153846153846156cm,y=1.277777777777778cm]\draw[->,gray] (-3.5,0)--(9.5,0) node[right]{$x$};\draw[->,gray] (0,-1.7)--(0,1.9) node[above]{$y$};\draw[blue,thick,solid,domain=-3:9,samples=180] plot(\x,{cos(deg(pi*(\x+1)/4))});\fill[orange!85!black] (0,0.7071067811865476) circle(2.2pt);\fill[orange!85!black] (-2,0.7071067811865476) circle(2.2pt);\fill[orange!85!black] (8,0.7071067811865476) circle(2.2pt);\foreach \a in {-1,3}{\draw[violet,dashed] (\a,-1.3)--(\a,1.3);\node[above] at (\a,1.35){$x=\a$};}\end{scope}\node[anchor=west,align=left,text width=5.0cm] at (8.4,4.0) {$f=\cos\frac{\pi(x+1)}4$\\[9pt]轴 $x=-1,3$\\[9pt]$0\to-2\to8$\\[9pt]三点高度 $\frac{\sqrt2}2$\\[9pt]周期可取 $2\times4=8$};
\node[text width=12.6cm,align=center] at (7.1,.65) {两次反射的合成位移，是两轴间距的两倍；这里只证明一个周期。};
\end{tikzpicture}\end{center}
